\begin{afterword}
\summaryhead

The introduction to the sixth and last book, \textsc{Becoming Stronger}, calls the book a call to action and points the reader to further reading. The book's definition of normative rationality, by probability theory and decision theory, is said to be standard in cognitive science. Its philosophy is more controversial: Yudkowsky one-boxes on Newcomb's problem, a minority view among decision theorists, and wants a rule for choosing priors, as Jaynes and Hutter do. Drescher's \textsc{Good and Real} reaches many of the same conclusions.

The introduction separates the philosophical dispute between Bayesians and frequentists from the practical choice of statistical methods, both of which have uses, and cites evidence that statistical training improves everyday reasoning while logic classes do not.

It then previews the book's three sequences: Yudkowsky's intellectual history, what it takes to solve a hard problem, and rationality groups. It lists open questions, credits the original posts with inspiring \textsc{Less Wrong}, effective altruism and the Center for Applied Rationality, and ends by calling the art of rationality young.

\responsehead

The introduction is a reading list and a preview, and it is judged as one. Where it describes positions and sources, the descriptions are accurate wherever I could check them.

The claim that one-boxing is a minority view among decision theorists is supported by both PhilPapers surveys. In 2009, one box drew 21.3 per cent of all respondents against 31.4 per cent for two boxes. In 2020, it drew 22.7 per cent of decision theorists' answers other than ``Other.'' The description of Talbott's encyclopedia entry is accurate: it names Jaynes among the ``Objective Bayesians,'' who hold that priors are rationally constrained. The reply that frequentism is the more subjective view is the stopping-rule argument, and the introduction presents it as a reply by two named people. The frequentist answer to it is discussed in the notes on ``Beautiful Probability.''

The introduction also says plainly that the book's philosophy is ``more controversial'' than its definitions, and it lists testing the book's own prescriptions among the tasks for the next generation of primers.

One claim I could not check. The introduction says that training in statistics, and in some other fields such as psychology, improves reasoning outside the classroom, while classes in logic and informal fallacies have not proven as useful. The two studies it cites are, by their titles and the summaries available to me, studies of statistical training. I could not read them in full, so I could not verify the parts about psychology and about logic classes.

The rest is what introductions do: a preview of the sequences, whose names and order match the book, and an account of the book's influence. Its statement of why the essays were written agrees with Yudkowsky's own words two posts later, that the writing was meant ``to leave a trail to where I ended up by accident.''

In short: a reading list and preview whose descriptions of positions and sources are accurate as far as they could be checked.
\end{afterword}
